% !TeX root = ../../main.tex
\subsection{1次関数}

\begin{definitionbox}{\textbf{1次関数}}
    \label{def:一次関数}
    $y$ が $x$ の関数で，$y$ が $x$ の1次式で表されるとき、\textbf{$y$ は $x$ の1次関数である}という.
    一般に、1次関数は次のように表される。
    \[\fitblankbf[]{y = ax + b}^{*} \qquad (a, b \text{は定数},\quad a \neq 0)\]

    \vspace{0.5em}
    このとき、$a$ を\fitblankbf[]{傾き} あるいは \fitblankbf[]{変化の割合}, $b$ を\fitblankbf[]{切片} という.

    \vspace{1.5em}
    また、$a = \dfrac{\fitblankbox{y \text{の増加量}}}{\fitblankbox{x \text{の増加量}}}$ で求めることができる。
\end{definitionbox}
\vspace{-2.5em}

% 例題3.1
\begin{example}
    グラフの傾きが $3$ で，点 $(-1,\ 5)$ を通る直線の式を求めなさい。
\end{example}

\vspace{-1.5em}
\begin{proof}\mbox{}\\
    傾きが $3$ であるから，直線の式は次のように表すことができる。
    \[y = 3x + b\]

    \begin{hiddenanswer}
        $x = -1$, $y = 5$ をこの式に代入すると
        \begin{align*}
            5 &= 3 \times (-1) + b \\
            b &= 8
        \end{align*}
        よって $y = 3x + 8$
        \finalanswer{答え：$y = 3x + 8$}
    \end{hiddenanswer}
\end{proof}
\vspace{-2em}

% 例題3.2
\begin{example}
    $2$ 点 $(-5,\ 6)$, $(3,\ -2)$ を通る直線の式を求めなさい。
\end{example}

\begin{proof}\mbox{}\\
    直線の傾きは $\dfrac{-2 - 6}{3 - (-5)} = -1$

    \vspace{0.5em}
    よって，直線の式は $y = -x + b$

    \begin{hiddenanswer}
        $x = -5$, $y = 6$ をこの式に代入すると
        \begin{align*}
            6 &= -1 \times (-5) + b \\
            b &= 1
        \end{align*}
        よって $y = -x + 1$
        \finalanswer{答え：$y = -x + 1$}
    \end{hiddenanswer}
\end{proof}

\vfill
\noindent\rule{\textwidth}{0.4pt}

\vspace{0.2em}
\noindent
$^{*}$ $b = 0$ のとき $\bm{y = ax}$ となり，$y$ は $x$ に比例する.すなわち、比例は1次関数の \uwave{特別な場合} である.

\newpage

\begin{theorembox}{\textbf{1次関数のグラフ}}
    一次関数：$y = ax + b$ のグラフは \fitblankbf[]{切片} を通る \fitblankbf[]{直線}である.
    \begin{enumerate}
        \item $a > 0$ のとき，グラフは\fitblankbf[]{右上がり}である.
        \item $a < 0$ のとき，グラフは\fitblankbf[]{右下がり}である.
    \end{enumerate}
\end{theorembox}

\vspace{-2.5em}
% 例題3.3
\begin{example}
    グラフが下の図のような直線になる $1$ 次関数の式をそれぞれ求めなさい。

    \vspace{-0.3em}
    \begin{center}
    \begin{tikzpicture}[scale=0.36, x=1cm, y=1cm]
        \tikzset{
            gline/.style={semithick},
        }

        % 格子
        \draw[step=1, gray!40, very thin] (-7,-7) grid (7,7);

        % 座標軸
        \draw[thick] (-7,0) -- (7,0) node[right] {$x$};
        \draw[thick] (0,-7) -- (0,7) node[above] {$y$};
        \node[anchor=north east] at (-0.15,-0.15) {$O$};
        \node[anchor=north] at (5,0) {$5$};
        \node[anchor=north] at (-5,0) {$-5$};
        \node[anchor=west] at (0,5) {$5$};
        \node[anchor=east] at (0,-5) {$-5$};

        % (1) y = -3x + 3
        \draw[gline] ({-4/3},7) -- ({10/3},-7);
        \draw[dotted] (1,0) -- (1,3) -- (0,3);
        \fill (1,0) circle (1.4pt);
        \fill (0,3) circle (1.4pt);
        \node[anchor=west] at (3.5,-6.5) {{\small（1）}};

        % (2) y = x + 2
        \draw[gline] (-7,-5) -- (5,7);
        \draw[dotted] (-2,0) -- (-2,2) -- (0,2);
        \fill (-2,0) circle (1.4pt);
        \fill (0,2) circle (1.4pt);
        \node[anchor=west] at (5,6.5) {{\small（2）}};
    \end{tikzpicture}
    \end{center}
    \vspace{-0.5em}
\end{example}

\vspace{-2em}
\begin{proof}\mbox{}\\
    \begin{hiddenanswer}
        （1）点 $(0,\ 3)$ を通るから，切片は $3$

        点 $(1,\ 0)$ も通るから，傾きは $\dfrac{0 - 3}{1 - 0} = \dfrac{-3}{1} = -3$

        よって $y = -3x + 3$
        \finalanswer{（1）$y = -3x + 3$}

        （2）割愛
        \finalanswer{（2）$y = x + 2$}
    \end{hiddenanswer}
\end{proof}

\vspace{-2em}
% 例題3.4
\begin{example}
    直線 $l$, $m$ が，それぞれ次の式で表されるとき，$l$, $m$ の交点の座標を求めなさい。

        \[l: y = 3x + 5,\ m: y = -x + 3\]
\end{example}

\vspace{-2em}
\begin{proof}\mbox{}\\
    \begin{hiddenanswer}
        \[
        \begin{cases}
            y = 3x + 5 \\
            y = -x + 3
        \end{cases}
        \]
        を解くと
        \[x = -\dfrac{1}{2},\ y = \dfrac{7}{2}\]

        よって，$l$, $m$ の交点の座標は $\left(-\dfrac{1}{2},\ \dfrac{7}{2}\right)$
        \finalanswer{答え：$\left(-\dfrac{1}{2},\ \dfrac{7}{2}\right)$}
    \end{hiddenanswer}
\end{proof}
\newpage

% 練習問題3.1
\begin{exercise}
    \noindent
    \begin{tabular}[t]{@{}p{0.50\linewidth}@{\hspace{1.2em}}p{0.44\linewidth}@{}}
    \vspace{0pt}
    右の図のように，直線 $y = \dfrac{1}{3}x + 4$ が $y$ 軸と点 A で交わっている。この直線と直線 $y = x$ の交点を P とする。
    \vspace{0.5em}

    \begin{enumerate}[label=(\arabic*), leftmargin=1.8em, itemsep=0.3em, topsep=0.2em, parsep=0pt]
        \item 点 P の座標を求めなさい。
        \item 点 O を通り，$\triangle OAP$ の面積を $2$ 等分する直線の式を求めなさい。
    \end{enumerate}
    &
    \begin{minipage}[t]{\linewidth}
    \vspace{0pt}
    \centering
    \begin{tikzpicture}[scale=0.48, x=1.25cm, y=1.25cm]
        \tikzset{
            gline/.style={semithick},
        }

        % 格子
        \draw[step=1, gray!40, very thin] (0,0) grid (8,8);

        % 座標軸
        \draw[thick, ->] (0,0) -- (8.5,0) node[right] {$x$};
        \draw[thick, ->] (0,0) -- (0,8.5) node[above] {$y$};
        \node[anchor=north east] at (-0.12,-0.12) {$O$};

        % △OAP
        \fill[pattern=dots, pattern color=gray!55] (0,0) -- (0,4) -- (6,6) -- cycle;

        % 直線
        \draw[gline] (0,4) -- (8,{8/3+4});
        \draw[gline] (0,0) -- (8,8);

        % 点・ラベル
        \node[anchor=east] at (-0.15,4) {$A$};
        \node[anchor=south west] at (5.5,5) {$P$};
        \node at (3.5,2) {$y = x$};
        \node[anchor=west] at (0.8,6) {$y = \dfrac{1}{3}x + 4$};

        % 交点・切片（最前面）
        \fill (0,4) circle (2pt);
        \fill (6,6) circle (2pt);
    \end{tikzpicture}
    \end{minipage}
    \end{tabular}
\end{exercise}

\begin{proof}\mbox{}\\
    \begin{hiddenanswer}
        （1）連立方程式
        \[
        \begin{cases}
            y = \dfrac{1}{3}x + 4 \\
            y = x
        \end{cases}
        \]
        を解くと $x = 6$, $y = 6$

        よって，点 $P$ の座標は $(6,\ 6)$
        \finalanswer{（1）$(6,\ 6)$}

        （2）点 O を通り，$\triangle OAP$ の面積を $2$ 等分する直線は，線分 AP の中点を通る。

        線分 AP の中点を $M$ とすると，$M$ の座標は
        \[\dfrac{0 + 6}{2} = 3,\ \dfrac{4 + 6}{2} = 5\]
        より $M$ の座標は $(3,\ 5)$

        よって，直線 OM の傾きは $\dfrac{5}{3}$ であるから，求める直線の式は $y = \dfrac{5}{3}x$
        \finalanswer{（2）$y = \dfrac{5}{3}x$}
    \end{hiddenanswer}
\end{proof}

\newpage

% 練習問題3.2
\begin{exercise}
    \noindent
    \begin{tabular}[t]{@{}p{0.50\linewidth}@{\hspace{1em}}p{0.44\linewidth}@{}}
    \vspace{0pt}
    $\mathrm{AB} = 3\ \mathrm{cm}$，$\mathrm{BC} = 4\ \mathrm{cm}$ である長方形 $\mathrm{ABCD}$ の周上を，頂点 A から毎秒 $1\ \mathrm{cm}$ の速さで，点 D を通り点 C まで動く点 P がある。P が頂点 A を出発してから $x$ 秒後の四角形 $\mathrm{ABCP}$ の面積を $y\ \mathrm{cm}^2$ とするとき，次の問いに答えなさい。
    \vspace{0.5em}

    \begin{enumerate}[label=(\arabic*), leftmargin=1.8em, itemsep=0.35em, topsep=0.2em, parsep=0pt]
        \item $0 \leqq x \leqq 4$ のとき，$y$ を $x$ の式で表しなさい。
        \item $4 \leqq x \leqq 7$ のとき，$y$ を $x$ の式で表しなさい。
        \item $0 \leqq x \leqq 7$ のときの $x$ と $y$ の関係をグラフに表しなさい。
        \item $y = 8$ となる $x$ の値をすべて求めなさい。
    \end{enumerate}
    &
    \begin{minipage}[t]{\linewidth}
    \vspace{0pt}
    \centering
    % 長方形の図
    \begin{tikzpicture}[scale=0.55, x=1cm, y=1cm]
        \coordinate (A) at (0,3);
        \coordinate (B) at (0,0);
        \coordinate (C) at (4,0);
        \coordinate (D) at (4,3);
        \coordinate (P) at (1.2,3);

        \fill[pattern=dots, pattern color=gray!55] (A) -- (B) -- (C) -- (P) -- cycle;
        \draw[semithick] (A) -- (B) -- (C) -- (D) -- cycle;
        \draw[semithick, -{Stealth[length=2mm]}] (P) -- (2.6,3);

        \node[anchor=south east] at (A) {A};
        \node[anchor=north east] at (B) {B};
        \node[anchor=north west] at (C) {C};
        \node[anchor=south west] at (D) {D};
        \node[anchor=south] at (P) {P};

        \draw[densely dashed] (-0.35,0) arc (180:90:0.35);
        \node[anchor=east] at (-0.45,1.5) {$3\ \mathrm{cm}$};
        \draw[densely dashed] (0,-0.35) arc (-90:0:0.35);
        \node[anchor=north] at (2,-0.45) {$4\ \mathrm{cm}$};
    \end{tikzpicture}

    \vspace{0.8em}

    % グラフ用の空欄
    \begin{tikzpicture}[scale=0.38, x=1cm, y=1cm]
        \draw[step=1, gray!40, very thin] (0,0) grid (8,14);
        \draw[thick, ->] (0,0) -- (8.5,0) node[right] {$x$ (秒)};
        \draw[thick, ->] (0,0) -- (0,14.5) node[above] {$y$ ($\mathrm{cm}^2$)};
        \node[anchor=north east] at (-0.12,-0.12) {O};
    \end{tikzpicture}
    \end{minipage}
    \end{tabular}
\end{exercise}

\begin{proof}\mbox{}\\
    \begin{hiddenanswer}
        （1）$0 \leqq x \leqq 4$ のとき

        四角形 ABCP は $\mathrm{AP} \parallel \mathrm{BC}$ の台形で，$\mathrm{AP} = x\ \mathrm{cm}$ であるから，その面積は
        \[\dfrac{1}{2} \times (x + 4) \times 3 = \dfrac{3}{2}x + 6\ (\mathrm{cm}^2)\]
        よって $y = \dfrac{3}{2}x + 6$
        \finalanswer{（1）$y = \dfrac{3}{2}x + 6$}

        （2）$4 \leqq x \leqq 7$ のとき

        四角形 ABCP は $\mathrm{AB} \parallel \mathrm{PC}$ の台形で，$\mathrm{PC} = 7 - x\ (\mathrm{cm})$ であるから，その面積は
        \[\dfrac{1}{2} \times \{(7 - x) + 3\} \times 4 = -2x + 20\ (\mathrm{cm}^2)\]
        よって $y = -2x + 20$
        \finalanswer{（2）$y = -2x + 20$}

        （3）（1），（2）より，$x$ と $y$ の関係をグラフに表すと，次のようになる。

        \begin{center}
        \begin{tikzpicture}[scale=0.42, x=1cm, y=1cm]
            \draw[step=1, gray!40, very thin] (0,0) grid (8,14);
            \draw[thick, ->] (0,0) -- (8.5,0) node[right] {$x$ (秒)};
            \draw[thick, ->] (0,0) -- (0,14.5) node[above] {$y$ ($\mathrm{cm}^2$)};
            \node[anchor=north east] at (-0.12,-0.12) {O};
            \node[anchor=east] at (0,6) {$6$};
            \node[anchor=east] at (0,12) {$12$};
            \node[anchor=north] at (4,0) {$4$};
            \node[anchor=north] at (7,0) {$7$};

            \draw[semithick] (0,6) -- (4,12) -- (7,6);
            \draw[dotted] (4,0) -- (4,12) -- (0,12);
            \draw[dotted] (7,0) -- (7,6);
            \fill (0,6) circle (1.6pt);
            \fill (4,12) circle (1.6pt);
            \fill (7,6) circle (1.6pt);
        \end{tikzpicture}
        \end{center}

        （4）グラフより，$y = 8$ となるのは $0 \leqq x \leqq 4$ のときに $1$ 回，$4 \leqq x \leqq 7$ のときに $1$ 回あることがわかる。

        $0 \leqq x \leqq 4$ のとき，$y = 8$ を $y = \dfrac{3}{2}x + 6$ に代入すると
        \[8 = \dfrac{3}{2}x + 6\]
        これを解くと $x = \dfrac{4}{3}$

        $4 \leqq x \leqq 7$ のとき，$y = 8$ を $y = -2x + 20$ に代入すると
        \[8 = -2x + 20\]
        これを解くと $x = 6$

        したがって，求める $x$ の値は $x = \dfrac{4}{3},\ 6$
        \finalanswer{（4）$x = \dfrac{4}{3},\ 6$}
    \end{hiddenanswer}
\end{proof}
